\documentclass[11pt]{article}
\usepackage[utf8]{inputenc}
\usepackage[T1]{fontenc}
\usepackage{amsmath,amssymb}
\usepackage{graphicx}
\usepackage{booktabs}
\usepackage{array}
\usepackage[margin=1in]{geometry}
\usepackage{hyperref}

\title{La idea de Francisco: super-nodos PWT y restriccion de $W$\\
\large (Notebook: 06\_pwt\_marmolejo)}
\author{}
\date{}

\begin{document}
\maketitle

\begin{abstract}
Francisco Marmolejo propuso una manera de re-formular el problema de
``elegir el siguiente pool'' en Greedy Dynamic, agregando un super-nodo
$w_T$ por cada subconjunto $T$ de la poblacion ya cubierta $S$. Si
existe una $\alpha$-aproximacion para el problema de un solo paso (PWT),
componiendola debiera dar una garantia $(1 - e^{-\alpha})$ para Greedy
Dynamic.

\textbf{Resumen:} la formulacion escalar (peso, prob, util) reproduce
exacto al greedy myopic con la lectura ``all-clear'', pero \emph{no}
cierra el lookahead gap, porque es ciega a la PMF de conteo. Resolvemos
empiricamente el problema de complejidad: la heuristica \texttt{partner}
permite reducir $W$ de tamano $2^{|S|}$ a un top-$L$ con $L \approx 1$--$3$
en promedio en seis regimenes diversos. Lo que queda abierto: probar
submodularidad adaptativa de $g_h(T,U)$ con observacion de conteo.
\end{abstract}

%% ====================================================================
\section{Setup}

Poblacion $N$, presupuesto $B$, tamano maximo de pool $G$, prior
$p \in [0,1]^N$, utilidad $u(F, Z) = \sum_{i \text{ proven healthy}} u_i$.
El test es \emph{aumentado}: $r(t, Z) = |t \cap Z|$, el conteo exacto.

Despues de $k$ pruebas con pools $t_1, \dots, t_k$, sea
$S = \bigcup_{j} t_j$ y $V = N \setminus S$. La idea de Francisco
agrega super-nodos $W = \{w_T : T \subseteq S, T \ne \emptyset\}$ con

\[
\text{peso}(w_T) = |T|, \quad
\text{prob}(w_T) = \prod_{i \in T}(1 - p_i), \quad
\text{util}(w_T) = \sum_{i \in T} u_i.
\]

\textbf{PWT (problema de un paso):} elegir $t \subseteq V \cup W$ con
$\sum \text{pesos} \le G$ y $|t \cap W| \le 1$, maximizando la utilidad
esperada.

\textbf{Algoritmo:} para cada $w_T \in W$ (incluyendo el caso ``ningun
super-nodo''), resolver/aproximar el sub-problema y quedarse con el
mejor.

%% ====================================================================
\section{Formulacion escalar es exacta para myopic}

\textbf{Claim.} La formulacion PWT con $\text{prob}(w_T) = \prod(1-p_i)$
reproduce \emph{exactamente} el valor de greedy myopic.

\textbf{Por que.} El score myopic estandar de un pool $t$ es
\[
\big(\textstyle\sum_{i \in t} u_i\big) \cdot
\textstyle\prod_{i \in t}(1 - p_i).
\]
Particionando $t = U \cup T$ con $U \subseteq V$ y $T \subseteq S$:
\[
\big(\textstyle\sum_{i \in U} u_i + \textstyle\sum_{i \in T} u_i\big) \cdot
\textstyle\prod_{i \in U}(1 - p_i) \cdot
\textstyle\prod_{i \in T}(1 - p_i)
=
(\text{util}_U + \text{util}_T) \cdot \text{prob}_U \cdot \text{prob}_T.
\]
Sumativo en utilidades, multiplicativo en probabilidades. Los factores
de $T$ entran exactamente como los de $U$.

\textbf{Verificado empiricamente} sobre $n=6$, $G=3$, $B=2$ con prior
y posterior arbitrarios: el pool elegido y el valor coinciden
bit-a-bit con la enumeracion completa.

\textbf{Lectura alternativa (OR-event).} Codex propuso
$\text{prob}(w_T) = 1 - \prod(1-p_i)$ — esto resuelve un \emph{objetivo
distinto} y elige peor pool. Es la lectura que rompe; la lectura
all-clear es la correcta.

%% ====================================================================
\section{Defecto: dos $T$ con mismo escalar pero distinto valor de info}

Aunque la formulacion escalar es exacta para el paso myopic, es
\emph{ciega a la PMF de conteo}. Esto importa cuando hay lookahead.

\textbf{Construccion.} $T_1 = \{a\}$ con $p_a = 0.5$, y $T_2 = \{b, c\}$
con $p_b = p_c = 1 - 1/\sqrt{2} \approx 0.293$. Tienen identicos:
\begin{itemize}
\item OR-prob: $\Pr(Z \cap T \ne \emptyset) = 0.5$
\item All-clear prob: $\prod(1 - p_i) = 0.5$
\item util si $u_a = u_b + u_c$
\end{itemize}
Pero la PMF del conteo $r$ difiere:
\[
\Pr_{T_1}(r) = (0.50, \, 0.50), \quad
\Pr_{T_2}(r) = (0.50, \, 0.41, \, 0.09).
\]
La entropia de Shannon:
$H(r_{T_1}) = 1.00$ bits, $H(r_{T_2}) = 1.33$ bits. $T_2$ acarrea
estrictamente mas informacion sobre $Z$.

\textbf{Implicacion.} Cualquier algoritmo de un paso que use solo el
escalar $(|T|, \text{prob}, \text{util})$ trata a $T_1$ y $T_2$
identico, pero un paso con lookahead que actualice por el resultado
$r > 0$ los va a distinguir. Por lo tanto la formulacion escalar de
PWT \emph{no} cierra el lookahead gap.

%% ====================================================================
\section{El lookahead gap crece con $B$}

Para la instancia $n=6$, $G=3$, $p$ y $u$ fijos:

\begin{table}[h]
\centering
\begin{tabular}{rrrr}
\toprule
$B$ & greedy(true) & DP optimum & gap \\
\midrule
2 & $18.4414$ & $18.5817$ & $+0.1403$ \\
3 & $20.4740$ & $22.4339$ & $+1.9598$ \\
\bottomrule
\end{tabular}
\caption{Lookahead gap en funcion del presupuesto. greedy(true) se
obtiene simulando la politica greedy sobre cada $z$ y promediando.}
\end{table}

PWT escalar reproduce exacto a greedy myopic, asi que vive en la
columna ``greedy(true)''. \textbf{No reduce el gap.}

%% ====================================================================
\section{Q3: restringir $W$ a top-$L$}

$|W| = 2^{|S|} - 1$ es exponencial. Para hacer la enumeracion practica
queremos quedarnos con un top-$L$ por una heuristica barata. Definimos
$L_{\min}$ = el menor $L$ tal que el top-$L$ contiene al $T^*$ optimo.

\textbf{Heuristicas probadas.}
\begin{itemize}
\item \texttt{self}: $\big(\sum u_T\big) \cdot \prod(1 - p_T)$
  — valor del pool $= T$ solo.
\item \texttt{prob}: $\prod(1 - p_T)$.
\item \texttt{util}: $\sum u_T$.
\item \texttt{ent\_$\lambda$}: $\texttt{self}(T) + \lambda \cdot H(r_T)$
  — aumentado con entropia.
\item \texttt{partner}: $(\sum u_T + u^*) \cdot \prod(1 - p_T) \cdot p^*$
  donde $(u^*, p^*)$ es la utilidad/prob del \emph{globalmente mejor
  pool de $V$ de tamano $G - |T|$}, precomputado.
\end{itemize}

\textbf{Resultado del sweep} (K=20 trials por regimen, $|W|$ es el
mean, format mean/max):

\begin{table}[h]
\centering
\begin{tabular}{l r r r r}
\toprule
Regimen & $|W|$ & partner & self & prob \\
\midrule
baseline ($n=10$, $G=4$, $p$ baja) & $79.7$ & $\mathbf{1.6}/13$ & $4.3/21$ & $16.1/98$ \\
prevalencia alta ($p \sim U[.4,.7]$) & $79.7$ & $2.8/14$ & $\mathbf{0.7}/2$ & $4.8/42$ \\
$n$ grande ($n=15$, $G=5$) & $278.3$ & $\mathbf{1.1}/4$ & $12.7/132$ & $38.3/233$ \\
historia profunda (3 pasos) & $62.9$ & $\mathbf{1.1}/3$ & $3.8/18$ & $14.1/40$ \\
prevalencia bimodal & $79.7$ & $\mathbf{1.0}/1$ & $8.4/35$ & $22.5/86$ \\
utilidad outlier ($u_0 = 50$) & $79.7$ & $\mathbf{1.4}/12$ & $2.4/13$ & $12.8/82$ \\
\bottomrule
\end{tabular}
\caption{$L_{\min}$ por regimen. \texttt{partner} domina en 4 de 6
regimenes y se mantiene $\le 3$ en mean siempre. Importante: para
$n=15$ con $|W| = 278$, mean $L_{\min}$ se mantiene en $1.1$ — no
escala con $|S|$.}
\end{table}

\textbf{Por que \texttt{partner} funciona.} Es esencialmente una cota
inferior sobre $\text{val}(T)$ usando como sucedaneo el mejor pool
global de $V$. Es ajustada cuando ese mejor pool no cambia mucho con
el tamano de $T$ — lo cual ocurre en la mayoria de los regimenes.
Crucialmente, \emph{toma en cuenta el budget consumido por} $T$, cosa
que \texttt{self} ignora.

\textbf{Entropia no ayuda.} \texttt{ent\_$\lambda$} con $\lambda > 0$
empata o empeora a \texttt{self}. Razon: $T$'s con alta entropia son
tipicamente $T$'s grandes con probabilidades intermedias, que son
malos pools myopicos. La informacion de pasos futuros no es la senal
correcta para rankear el paso myopic.

%% ====================================================================
\section{Caso adversarial}

Construimos $S = \{0,1,2,3\}$, $V = \{4,5\}$, $G = 4$,
$p = (.01, .01, .60, .05, .05, .50)$,
$u = (1, 1, 100, 50, 60, 80)$.

Pool optimo: $\{3, 4\}$ con $\text{val} = 99.27$.
$T^* = \{3\}$ (singleton). \texttt{self} ranquea $\{2, 3\}$ arriba
porque $\text{util}_T = 150$ domina pese a $\text{prob}_T = 0.4$. Pero
$\text{val}(\{2,3\}) = 75.81$ — el factor $0.4$ multiplica con el
factor de $V$ y hunde el conjunto.

\begin{table}[h]
\centering
\begin{tabular}{l r r}
\toprule
Heuristica & $L_{\min}$ & Fraccion de $|W|$ \\
\midrule
\texttt{partner} & $\mathbf{2}$ & $13.3\%$ \\
\texttt{prob} & $4$ & $26.7\%$ \\
\texttt{ent\_1} & $8$ & $53.3\%$ \\
\texttt{self} & $8$ & $53.3\%$ \\
\texttt{util} & $12$ & $80.0\%$ \\
\texttt{rand} & $15$ & $100.0\%$ \\
\bottomrule
\end{tabular}
\caption{Adversarial worst-case. $|W| = 15$. \texttt{partner} mantiene
$L_{\min}$ chico (top-2), \texttt{self} requiere bajar a la posicion 8
de 15 — \emph{53\%} de la lista.}
\end{table}

%% ====================================================================
\section{Conclusiones y trabajo abierto}

\textbf{Lo que la formulacion PWT da:}
\begin{enumerate}
\item Equivalencia myopic exacta con la lectura all-clear. El escalar
  $(|T|, \prod(1-p), \sum u)$ basta para reproducir la decision myopic.
\item Restriccion practica: la heuristica \texttt{partner} reduce
  $|W| = 2^{|S|}$ a top-$L$ con $L \sim 1$--$3$ en promedio en
  $n = 15$ con $|W| = 278$.
\end{enumerate}

\textbf{Lo que falta:}
\begin{enumerate}
\item \textbf{Cota worst-case sobre $L_{\min}$ con \texttt{partner}}.
  Empiricamente max $L_{\min} \le 14$ en todos los regimenes, pero no
  hay prueba teorica ni adversarial conocido.
\item \textbf{Submodularidad adaptativa de $g_h(T, U)$ con observacion
  de conteo}. Sin esta, una $\alpha$-aproximacion por paso \emph{no}
  compone a una garantia global $(1 - e^{-\alpha})$. Esta es la
  pregunta Q4 en \texttt{docs/notes/pwt\_submodularity.md}.
\item \textbf{Ranking lookahead-aware}. La entropia simple no funciona,
  pero un ranking que use la \emph{ganancia esperada de informacion}
  podria preservar garantias multi-paso.
\end{enumerate}

\textbf{Veredicto practico.} PWT con \texttt{partner} y $L = 5$ es una
implementacion valida del paso de Greedy Dynamic, equivalente al
greedy actual pero mas modular para analisis teorico. \textbf{Veredicto
teorico.} Sin Q4, la formulacion es una re-organizacion limpia, no una
reduccion garantizada.

\end{document}
